
@@M0@@,
évidemment unique.

\textbf{Proposition 4.8.} \uline{Soit} @@M1@@ \uline{un entier} @@M2@@~; \uline{dans le} @@M3@@-\uline{anneau libre engendré} \uline{sur} @@M4@@ \uline{par une indéterminée} @@M5@@, \uline{l'idéal engendré par les} @@M6@@, @@M7@@, \uline{est} \uline{un} @@M8@@-\uline{idéal}.

Posons @@M9@@, et soit @@M10@@ l'idéal engendré par les @@M11@@ pour @@M12@@~; posons @@M13@@. Il faut montrer que l'homomorphisme composé : @@M14@@ s'annule sur @@M15@@, i.e. que @@M16@@, @@M17@@, a pour image @@M18@@ dans @@M19@@. Or @@M20@@, et @@M21@@ est un @@M22@@-homomorphisme~; il suffit donc de montrer que pour tout anneau @@M23@@ et toute série @@M24@@, dont les coefficients sont nuls en degrés @@M25@@, @@M26@@ si @@M27@@. Mais il suffit pour prouver cette dernière propriété de la montrer pour @@M28@@, et @@M29@@. Plongeant @@M30@@ dans @@M31@@ grâce aux fonctions symétriques, on peut écrire : @@M32@@, d'où
@@M33@@
qui est de degré @@M34@@ en @@M35@@.

\textbf{Corollaire 4.9.} \uline{Le} @@M36@@-\uline{anneau engendré sur} @@M37@@ \uline{par le générateur} @@M38@@ \uline{soumis aux relations} @@M39@@ \uline{pour} @@M40@@, \uline{est isomorphe à} @@M41@@.

4.9.1. On aurait les mêmes énoncés en remplaçant les @@M42@@ par les @@M43@@. On peut également remplacer @@M44@@ par une famille d'indéterminées @@M45@@, soumises aux relations @@M46@@ pour @@M47@@, les @@M48@@ étant des entiers @@M49@@~; le @@M50@@-anneau obtenu est alors @@M51@@.

4.10. Soient @@M52@@ un anneau binômial, et @@M53@@ une @@M54@@-@@M55@@-algèbre augmentée. On désigne par @@M56@@ le @@M57@@-anneau libre engendré par @@M58@@ sur @@M59@@~; 
